% ============================================================
%  GnssProcessor: Mathematical Reference -- COMBINED VOLUME
%
%  This is the master document.  It stitches every component
%  document into one PDF via the `subfiles' package.  Each
%  component file (foundations.tex, slip_detection.tex, ...)
%  ALSO compiles on its own into a standalone per-component
%  PDF, sharing this file's preamble (gnssreport.sty).
%
%  Scope: equations, state vector layout, and design rationale
%  that Doxygen can't typeset cleanly.  Implementation-level
%  documentation (call signatures, parameter docs) remains in
%  the headers under include/ and is rendered via Doxygen.
%
%  Build:
%      latexmk -pdf main.tex            # the combined volume
%      latexmk -pdf slip_detection.tex  # one component, standalone
%  or, for every target at once:
%      ./build.ps1        (Windows)   /   ./build.sh   (Unix)
% ============================================================

\documentclass[11pt,letterpaper]{article}
\usepackage{gnssreport}     % shared preamble + notation (single source of truth)
\usepackage{subfiles}       % per-component PDFs + this combined volume from one tree
\addbibresource{\subfix{references.bib}}

\title{\bfseries GnssProcessor: Mathematical Reference}
\author{}
\date{Revision date: \today}

\begin{document}
\maketitle
\tableofcontents
\bigskip

\section{Introduction}\label{sec:intro}

This document is the mathematical companion to \texttt{GnssProcessor}, a
C\texttt{++}17 post-processing and real-time GNSS engine implementing
PPP, PPP-RTK (single-receiver and network), and double-differenced
baseline modes.  Its scope is the estimation mathematics and design
rationale: the material that ages slowly because it is settled theory
or algebra, not churning code.

The reference is organised as a set of \emph{component documents}, each of
which also compiles to a standalone PDF.  This combined volume gathers
them in roughly the order the engine applies them:
\begin{itemize}[leftmargin=*]
\item \textbf{Foundations} (\Cref{sec:notation,sec:state,sec:obs,sec:diff}):
      notation, the variable-dimension state vector, the
      \emph{uncombined} pseudorange / carrier-phase observation model, and
      the between-satellite single-difference (BSSD) and double-difference
      (DD) transformations.  Every other component builds on this.
\item \textbf{Reference systems and measurement models}
      (\Cref{sec:frames,sec:products,sec:atmos,sec:sitedisp,sec:obsphys}):
      the deterministic substrate and physics behind the observation
      model the Foundations document writes abstractly: the coordinate
      frames and time systems every other chapter assumes; the baseline
      satellite orbit/clock product evaluation that SSR refines;
      tropospheric and ionospheric delay modelling; the geophysical
      site displacements (solid Earth, ocean, and pole tides) that
      correct the receiver coordinate; and the antenna, attitude,
      relativity, and range corrections that map the raw measurement to
      the modelled observable.
\item \textbf{Estimation engine} (\Cref{sec:kf,sec:stoch,sec:estinfra}):
      the variable-dimension Kalman filter (Joseph-form update,
      prior-anchored iterated relinearisation, robust gating) and the
      backward RTS smoother; the measurement stochastic model and the
      a-posteriori quality-control layer (variance factor, Baarda
      data-snooping, minimal detectable bias); and the supporting
      infrastructure that bootstraps the filter: the weighted
      least-squares seed, the first-epoch residual-difference state
      initialisation, and the pivot selection / disturbance logic.
\item \textbf{Cycle-slip detection and repair} (\Cref{sec:slip}): the
      linear-combination test statistics (Melbourne--W\"ubbena,
      geometry-free, GRAPHIC, extra-wide-lane) and the TurboEdit
      $2\times2$ integer repair.
\item \textbf{Integer ambiguity resolution} (\Cref{sec:ar}): the
      \textsc{MLAMBDA} integer least-squares search, wide-lane /
      narrow-lane decomposition, validation, constrained feedback, and
      the (default-off) network-producer AR machinery.
\item \textbf{Inertial aiding} (\Cref{sec:imu}): the strapdown
      mechanisation and 15-state INS error model of the tightly-coupled
      IMU path.
\item \textbf{Real-time SSR corrections} (\Cref{sec:ssr}): how State
      Space Representation orbit, clock, and bias corrections refine the
      undifferenced observation model.
\item \textbf{Real-time engine} (\Cref{sec:rt}): the streaming
      architecture: NTRIP transport, RTCM/MSM decoding, observation/SSR
      time synchronisation, and the per-epoch real-time update path.
\item \textbf{Network mode} (\Cref{sec:net}): the multi-station
      PPP-RTK network estimator that \emph{generates} the SSR-style
      corrections consumed by the single-receiver modes: the joint
      observation model, the S-basis / minimum-constraint datum, the
      per-station Schur reduction, and the single-concern seam
      architecture of the per-epoch producer loop.
\item \textbf{Between-epoch carry} (\Cref{sec:net-carry}): how
      information reaches that arrowhead system from outside the
      current epoch: the three folds an arrowhead admits and where each
      must be applied, the double-difference cycle basis (the spanning
      tree's \emph{chords}) that is the only gauge-free quantity worth
      time-linking, the Woodbury algebra that folds it, and a fully
      numerical minimum-parameter example.
\item \textbf{Pipeline and architecture diagrams} (\Cref{sec:diagrams}):
      the per-mode data-flow figures that place every chapter above
      in the epoch loop.
\item \textbf{Generated reference tables} (\Cref{sec:gentables},
      appendix): processing modes, feature toggles, the cycle-slip
      detector registry, and the state-vector key schema, produced
      directly from the source tree by
      \texttt{scripts/gen\_report\_tables.py} rather than maintained by
      hand.
\end{itemize}

Implementation details (function signatures, class responsibilities,
config field reference) live in the headers under \texttt{include/}
and are rendered via Doxygen.  When this document references a
specific algorithmic component, it cites the responsible \emph{symbol}
(a function, class, or JSON config key) together with its source
file (e.g.\ \texttt{KalmanFilter::update} in
\texttt{src/core/filter/KalmanFilter.cpp}, or
\texttt{network.cold\_clock\_seed}).  Line numbers are deliberately
never cited: they rot within weeks, symbol names do not.

The notation follows the IGS conventions adopted by
\textcite{subirana2013gnss} and \textcite{teunissen2017handbook};
see \Cref{sec:notation} for the symbol table this document depends on.

% ------------------------------------------------------------
%  Component documents.  Each \subfile is also a standalone PDF.
% ------------------------------------------------------------
\part{Foundations}
\subfile{foundations}

\part{Reference systems and measurement models}
\subfile{frames_time}
\subfile{product_evaluation}
\subfile{atmosphere}
\subfile{site_displacement}
\subfile{observation_physics}

\part{Estimation engine}
\subfile{kalman_smoother}
\subfile{stochastic_qc}
\subfile{estimator_infrastructure}

\part{Cycle-slip detection and repair}
\subfile{slip_detection}

\part{Integer ambiguity resolution}
\subfile{ambiguity_resolution}

\part{Inertial aiding}
\subfile{imu_aiding}

\part{Real-time SSR corrections}
\subfile{ssr_corrections}

\part{Real-time engine}
\subfile{realtime_engine}

\part{Network mode}
\subfile{network}
\subfile{network_example}
\subfile{network_carry}

\part{Reference}
\subfile{pipeline_diagrams}

\appendix
\subfile{generated_tables}

\printbibliography

\end{document}
